\section{Method: alternative-dependent diagnosability}\label{sec:method}

\subsection{Definitions and terms}\label{sec:defs}

Let $\mathcal{A}$ be a set of design \emph{alternatives}. An alternative $a\in\mathcal{A}$ is an architecture together with the control structure it implies---for a generator, the machine topology and the converter control that the topology requires. Let $\mathcal{F}$ be a set of failure-mode classes; each class $f\in\mathcal{F}$ is parameterised by an \emph{extent} $e\in E_f$ expressed at the design level, so that the same $(f,e)$ can be realised in every alternative. For a winding fault the extent is the fraction of a winding branch that is short-circuited; for a bearing it might be a defect size; for a structure a crack length. The design-level extent does not fix the physical severity of the event (Section~\ref{sec:res-severity}), so the framing also names a \emph{physical severity axis} on which the alternatives can be compared. Let $\mathcal{S}$ be a set of candidate \emph{observation suites}; a suite $s\in\mathcal{S}$ is a set of measurements with a cost $c(s)$, and $\Phi(s)$ is the set of features computable from those measurements alone. Let $\mathcal{O}$ be the operating envelope, discretised into operating points $o$.

A \emph{trial} is a record of the system in alternative $a$ at operating point $o$ containing a reference interval $R$ in which the system is healthy and, immediately after it, a faulted interval $F$ in which failure mode $(f,e)$ is present. Trials may come from a physical test (test-stage instantiation) or from a validated simulation (design-stage instantiation, Section~\ref{sec:design-stage}). We use these terms throughout: alternative (not topology or machine, except when a specific machine is meant), trial, observation suite, extent (design level) and severity (physical).

\subsection{The signal-to-drift ratio}\label{sec:sdr}

The reference interval is split into two adjacent halves $R_1, R_2$, with $R_2$ adjacent to $F$. For a feature $\varphi\in\Phi(s)$ and a trial, let $\bar\varphi_{R_1}, \bar\varphi_{R_2}, \bar\varphi_F$ be the feature's mean values over the intervals, and define the relative change under fault, referenced to the half adjacent to the fault,
\begin{equation}
  \delta_\varphi = \frac{\bar\varphi_F - \bar\varphi_{R_2}}{|\bar\varphi_{R_2}|},
  \label{eq:delta}
\end{equation}
and the relative change between the two healthy halves,
\begin{equation}
  \delta^0_\varphi = \frac{\bar\varphi_{R_2} - \bar\varphi_{R_1}}{|\bar\varphi_{R_1}|},
  \label{eq:null}
\end{equation}
which measures how much the feature moves when nothing is wrong: sensor noise, drift of the operating point, controller activity, thermal effects. Because every trial has a healthy reference interval, $\delta^0_\varphi$ is available from every trial of alternative $a$, faulted or not, and its absolute value pooled over the $N_a$ trials forms an empirical null distribution of ordinary non-stationarity for that alternative and feature. The signal-to-drift ratio is
\begin{equation}
  \mathrm{SDR}_\varphi = \frac{|\delta_\varphi|}{q_\alpha\!\left(|\delta^0_\varphi|\right)},
  \label{eq:sdr}
\end{equation}
where $q_\alpha$ is the $\alpha$-quantile of the pooled null. With $\alpha=0.95$, $\mathrm{SDR}_\varphi\ge1$ means that the change under fault exceeds what the healthy system produces on its own in 95\,\% of comparable interval pairs.

\paragraph{What the SDR is statistically.} For a roughly Gaussian null, $q_{0.95}(|\delta^0|)\approx1.96\,\sigma$, so $\mathrm{SDR}\ge1$ is a two-sided 5\,\% test on the shift of interval means with a known change time: the fixed-sample mean-shift test, i.e.\ the known-onset special case of the generalised likelihood ratio of \citet{Basseville1993}. The SDR is therefore a trial-level standardised effect size (Cohen's $d$ on interval means, divided by 1.96), the detectable share defined below is an empirical power, and a share at a given $\alpha$ is one point of a receiver operating characteristic. Three assumptions are built in: the onset is known (the trial supplies it), the whole faulted interval is averaged, and the healthy reference lies immediately before the fault. The last assumption makes the null one of sub-second non-stationarity; a reference established at commissioning, months earlier, has a wider null, and Section~\ref{sec:res-robust} quantifies the difference on the demonstration data. Consequently the SDR characterises the simplest detector on the feature---a threshold on the relative change of interval means---and the design quantities derived from it are statements about that detector; a detector that combines features or accumulates evidence over time can detect below $\mathrm{SDR}=1$, which is why the oracle variant and the transfer experiment of step~6 are part of the method.

Four properties make the SDR suitable as a design measure. It is \emph{scale-free}: a current ratio, a voltage harmonic and a controller output are compared on one axis. It is \emph{calibrated} against each alternative's own non-stationarity, so a noisier alternative is not credited with detectability it does not have; the price is that the numerator and the denominator must be reported separately when alternatives are compared (Section~\ref{sec:res-where}), because a quieter null raises the SDR as much as a stronger signature does. It is \emph{per trial}, so its distribution over trials carries the dependence on operating point and extent. And it requires \emph{no detector}, which is what allows it to be computed before any algorithm is chosen. The relative form (\ref{eq:delta}) presumes a feature with a non-zero healthy level; for features whose healthy value is a noise floor an absolute-change variant with the same structure is used, and Section~\ref{sec:res-robust} shows that the two rank the features alike.

\paragraph{Safeguards.} Three safeguards are part of the definition. (i) \emph{Reliability screen.} On healthy trials---trials in which the ``faulted'' interval contains no fault, only the switching operation that would insert one---a feature should exceed $\mathrm{SDR}=1$ at the nominal rate $1-\alpha$. A feature whose count of such false alarms in the $n_h$ healthy trials is incompatible with that rate (one-sided binomial test at the 5\,\% level) is flagged as unreliable for that alternative and excluded from the design quantities; with $n_h=9$ and $\alpha=0.95$ this means three or more false alarms. Features whose null quantile exceeds one, i.e.\ whose healthy relative change exceeds 100\,\%, are excluded as uninformative. (ii) \emph{Fixed versus oracle feature.} A suite is evaluated by one \emph{fixed} feature---the feature with the highest median SDR over the trials of a failure-mode class among the screened features of the suite, chosen once---or by the \emph{oracle} feature, the best feature of each trial. The fixed variant is what a threshold detector achieves; the oracle variant is an upper bound for a detector that combines the suite's features, and because it takes a maximum over features its false-alarm rate on healthy trials is reported beside it. Since the fixed feature is chosen on the same trials it is evaluated on, a \emph{nested} variant chooses it on all but one tap pair and evaluates it on the held-out pair; both are reported. (iii) \emph{Uncertainty.} Every design quantity carries a bootstrap interval: over tap pairs (the physical fault locations) for statements about a failure-mode class, over operating points within an extent level for the minimum detectable extent, with Wilson intervals for shares.

\subsection{Derived design quantities and the definition}\label{sec:derived}

From the SDR distribution over trials we define, for each alternative $a$, suite $s$ and failure-mode class $f$:
\begin{itemize}
  \item the \emph{diagnosability map} $D_a(s,f,e,o)$: the SDR of the suite's fixed feature at extent $e$ and operating point $o$;
  \item the \emph{minimum detectable extent} $\mathrm{MDE}_a(s,f)$: the smallest tested extent level $e^\ast$ such that $\operatorname{median}_o D_a(s,f,e,o)\ge1$ for every tested $e\ge e^\ast$ (the limit-of-detection convention; equal extents realised at different winding locations form one level);
  \item the \emph{detectable share} $\mathrm{DS}_a(s,f\mid e\ge e_{\mathrm{req}})$: the fraction of faulted trials with $e\ge e_{\mathrm{req}}$ in which $\mathrm{SDR}\ge1$, a quantity written for a required extent so that it does not depend on how the test matrix samples the extent axis;
  \item the \emph{signature location}: the measurement group---stator currents, terminal voltages, controller internals, mechanical quantities, excitation---whose best screened feature attains the highest SDR in the largest share of trials.
\end{itemize}
\emph{Alternative-dependent diagnosability} is the map $(a, s, f)\mapsto(\mathrm{MDE}_a, \mathrm{DS}_a, \text{location})$. For a requirement $(f, e_{\mathrm{req}}, \alpha)$ the \emph{adequate} suites of alternative $a$ are those with $\mathrm{MDE}_a(s,f)\le e_{\mathrm{req}}$, and diagnosability is alternative-dependent for that requirement when the cheapest adequate suite differs between alternatives by more than the bootstrap uncertainty of the MDE. MDE and DS are the quantities against which requirements are written (``inter-turn faults of extent $\ge e_{\mathrm{req}}$ shall be detectable over the operating envelope at a per-trial false-alarm rate $\le1-\alpha$''); the signature location tells the designer \emph{why} an alternative behaves as it does.

\subsection{Design procedure}\label{sec:procedure}

\begin{figure}[t]
\centering
\resizebox{\textwidth}{!}{%
\begin{tikzpicture}[
  font=\small,
  box/.style={draw, rounded corners=2pt, align=left, inner sep=5pt, text width=3.9cm, minimum height=1.35cm},
  data/.style={draw, align=center, inner sep=4pt, text width=3.2cm, fill=black!4},
  arr/.style={-{Latex[length=2mm]}, thick},
  node distance=4mm and 6mm
]
\node[box] (s1) {\textbf{1 Frame}\\ alternatives $\mathcal{A}$, failure modes $\mathcal{F}$ with extent scale $E_f$, suites $\mathcal{S}$ with cost $c(s)$, envelope $\mathcal{O}$};
\node[box, right=of s1] (s2) {\textbf{2 Generate trials}\\ per $(a,f,e,o)$: reference interval $R$ + faulted interval $F$; simulation (design stage) or bench (test stage); healthy trials included};
\node[box, right=of s2] (s3) {\textbf{3 Compute}\\ features $\Phi(s)$ on windows; pooled null $q_\alpha(|\delta^0|)$ per $(a,\varphi)$; $\mathrm{SDR}=|\delta|/q_\alpha$; screen unreliable features};
\node[box, below=of s3] (s4) {\textbf{4 Map}\\ $D_a(s,f,e,o)$; $\mathrm{MDE}_a(s,f)$, $\mathrm{DS}_a(s,f)$; signature location; fixed and oracle variants};
\node[box, left=of s4] (s5) {\textbf{5 Select}\\ cheapest $s$ meeting the requirement per $a$; compare alternatives on $(c, \mathrm{MDE}, \mathrm{DS})$ with the other criteria};
\node[box, left=of s5] (s6) {\textbf{6 Assess transfer}\\ detector built on one alternative tested on the others under candidate referencing rules; fix the rule in the specification};
\draw[arr] (s1) -- (s2);
\draw[arr] (s2) -- (s3);
\draw[arr] (s3) -- (s4);
\draw[arr] (s4) -- (s5);
\draw[arr] (s5) -- (s6);
\node[data, below=6mm of s6] (out) {co-selected architecture + observation suite + referencing rule};
\draw[arr] (s6) -- (out);
\node[data, below=6mm of s5, text width=4.2cm] (req) {detectability requirement:\\ $\mathrm{MDE}\le e_{\mathrm{req}}$ at false-alarm $\le 1-\alpha$};
\draw[arr] (req) -- (s5);
\end{tikzpicture}}
\caption{The design procedure. Steps 3--4 are fully specified by Eqs.~(\ref{eq:delta})--(\ref{eq:sdr}); steps 1, 2, 5 and 6 are design decisions the method structures.}\label{fig:method}
\end{figure}


The procedure (Fig.~\ref{fig:method}) has six steps. Steps 3--4 are mechanical and fully specified by Eqs.~(\ref{eq:delta})--(\ref{eq:sdr}); steps 1, 2, 5 and 6 are design decisions, for which the method supplies the following default rules.
\begin{enumerate}
  \item \textbf{Frame.} Enumerate the alternatives $\mathcal{A}$, the failure-mode classes $\mathcal{F}$ with a design-level extent scale $E_f$ realisable in every alternative and a physical severity axis, the candidate observation suites $\mathcal{S}$, and the operating envelope $\mathcal{O}$. Set $e_{\mathrm{req}}$ as the smallest extent whose growth to an unacceptable consequence is faster than the inspection or maintenance interval, and $1-\alpha$ as the per-trial false-alarm rate the operator tolerates over that interval. Order the suites by cost; an ordinal rank by added hardware and integration is the default, a monetary cost including installation and maintenance is used when available.
  \item \textbf{Generate trials.} For each $(a,f,e,o)$ obtain trials with a healthy reference interval of at least twice the faulted duration, immediately followed by the faulted interval. Test at least the extent levels that bracket $e_{\mathrm{req}}$ and the corners of the operating envelope, and include enough healthy trials with the same timing for the screen (nine allow the binomial rule above). At the test stage the trials are prototype or benchmark experiments; at the design stage the faulted numerator is simulated and the null is supplied as described in Section~\ref{sec:design-stage}.
  \item \textbf{Compute.} Extract the features $\Phi(s)$ for all suites on sliding windows within the intervals, form the pooled null per alternative and feature, compute the SDR per trial, apply the screen.
  \item \textbf{Map.} Build $D_a(s,f,e,o)$ and summarise it into $\mathrm{MDE}_a$, $\mathrm{DS}_a$ and the signature location, for the fixed, nested and oracle variants, with bootstrap intervals.
  \item \textbf{Select.} Build the decision table: for each alternative, class and suite, and for the candidate values of $e_{\mathrm{req}}$ and $\alpha$, classify the suite as \emph{meets} when the bootstrap probability that $\mathrm{MDE}\le e_{\mathrm{req}}$ is at least 0.95, \emph{fails} when it is at most 0.05, and \emph{uncertain} otherwise. The adequate suite of an alternative is the cheapest that meets the requirement for \emph{every} failure-mode class in $\mathcal{F}$; an uncertain verdict is resolved by more trials at the extents around $e_{\mathrm{req}}$ or, failing that, by specifying the next suite as a fallback. The alternatives are then compared on the cost of their adequate suites together with the other selection criteria. The architecture decision is affected when, for some alternative, no affordable suite meets the requirement; otherwise the method changes the monitoring specification, not the architecture, and says so.
  \item \textbf{Assess transfer.} If the observation suite and its detector must serve a family of alternatives, test whether a detector built on one alternative's trials detects the others' faults under the candidate feature-referencing rules (physical units; scaled on the target's healthy trials; referenced to each trial's own preceding healthy interval). The rule that transfers becomes part of the monitoring specification.
\end{enumerate}

\subsection{Design-stage instantiation}\label{sec:design-stage}

The demonstration in this paper is a test-stage instantiation on built machines. A design-stage instantiation must supply the two ingredients of the SDR separately. The numerator---the change of a feature under a fault of extent $e$---is what a finite-element or circuit model of each alternative with the fault injected can provide \citep{Vaseghi2009}. The denominator---the healthy non-stationarity of sensor noise, controller activity, thermal drift and operating-point wander---is a property of hardware and environment that such a model does not reproduce, and a simulated SDR that supplied its own null would be optimistic by an unknown factor. The null must therefore come from measured healthy records of hardware of the same class (a predecessor product, a prototype of one alternative, or a benchmark such as the one used here), transferred to the other alternatives under the assumption that the healthy non-stationarity of a feature depends on the sensing and control chain more than on the machine; that assumption is itself testable with the method, by comparing the null tables of the alternatives for which hardware exists (Table~\ref{tab:null}). The validation of a simulated instantiation is then a comparison of simulated and measured SDR on a few physical trials, which is the role the three benches would play in a design study of the next generator.

\subsection{Scope}\label{sec:scope}

The method applies to failure-mode classes that can be seeded at a known extent, at a stationary operating point, with a healthy interval immediately before the fault, in every alternative; the extent scale must be realisable in all of them, and the physical severity axis is declared beside it. It evaluates detectability by a threshold on interval means, not the time to detection, and it stops at detection: isolation, severity estimation and remaining-life prediction presuppose detectability and can be built on the same trials. It does not model the cost of missed detection, which enters through $e_{\mathrm{req}}$ and $\alpha$ in the requirement. Run-to-failure data, in which the fault grows gradually and no healthy interval precedes a known onset, do not fit the trial structure and would need a different null.
